\documentclass[aspectratio=1610, 9pt]{beamer}
\usepackage{mobi-beamer}

\title[Pensamiento Sistémico e Introducció…]{Pensamiento Sistémico e Introducción a la Simulación}
\subtitle{}
\author{Alejandra Tabares Pozos}
\institute{Universidad de los Andes \\
 Departamento de Ingeniería Industrial}
\date{}

\begin{document}

%================================================
% SLIDE 1: PORTADA
%================================================
\begin{frame}
  \titlepage
\end{frame}

%================================================
% SLIDE 2: CONTEXTO
%================================================
\begin{frame}{Contexto del Curso}
\begin{itemize}
  \item \textbf{Curso:} Modelado Bajo Incertidumbre.
  \item \textbf{Rol en Ingeniería Industrial:} puente entre probabilidad (formal) y decisiones en sistemas reales (logística, manufactura, finanzas, energía).
  \item \textbf{Herramientas:}
    \begin{itemize}
      \item Teoría: variables aleatorias, procesos estocásticos, (nociones de) Markov.
      \item Práctica: simulación computacional en Python.
    \end{itemize}
  \item \textbf{Filosofía:} usar la computación para \emph{operacionalizar} el pensamiento probabilístico (no para reemplazarlo).
\end{itemize}
\end{frame}

%================================================
% SLIDE 3: OBJETIVOS
%================================================
\begin{frame}{Objetivos de Aprendizaje (Semana 1)}
Al finalizar la sesión, el estudiante podrá:
\begin{enumerate}
  \item Diferenciar formalmente entre sistemas determinísticos y estocásticos.
  \item Definir simulación y justificar cuándo es apropiada (y cuándo no).
  \item Identificar componentes del sistema: entidades, atributos, estado.
  \item Estimar numéricamente una métrica bajo incertidumbre con Monte Carlo (incluyendo error estadístico).
\end{enumerate}
\end{frame}

%================================================
% SLIDE 4: DEFINICIÓN DE SISTEMA
%================================================
\begin{frame}{Definición de Sistema}
\begin{block}{Definición (operativa)}
Un \textbf{sistema} es un conjunto de objetos que interactúan bajo reglas, con el fin de cumplir un propósito.
\end{block}

\vspace{0.3cm}
\textbf{Componentes clave:}
\begin{itemize}
  \item \textbf{Entidades:} objetos de interés (clientes, máquinas, vehículos, generadores).
  \item \textbf{Atributos:} propiedades de entidades (prioridad, velocidad, capacidad, estado de carga).
  \item \textbf{Actividades/eventos:} mecanismos de cambio (atender, fallar, producir, despachar).
\end{itemize}
\end{frame}

%================================================
% SLIDE 5: ESTADO
%================================================
\begin{frame}{El Estado del Sistema}
El \textbf{estado} es el conjunto mínimo de variables que permite describir el sistema en tiempo $t$ respecto al objetivo del estudio.

\vspace{0.3cm}
Sea $S(t)$ el vector de estado:
\[
S(t) = (x_1(t), x_2(t), \dots, x_n(t)).
\]

\textbf{Ejemplo (inventario):}
\[
S(t) = (\text{inventario disponible},\ \text{órdenes en tránsito},\ \text{tiempo hasta próxima revisión}).
\]
\end{frame}


\begin{frame}{Ejemplo de sistema 1: Microred con PV + batería + compra/venta a red}
\textbf{Propósito del estudio:} operación horaria para minimizar costo y riesgo de déficit.

\medskip
\textbf{Entidades (objetos del sistema):}
\begin{itemize}
  \item \textbf{Activos físicos:} Generación FV (PV), batería (BESS), carga agregada, punto de acople a red.
  \item \textbf{Agente de decisión:} EMS (controlador) que elige acciones $a_t$.
\end{itemize}

\textbf{Atributos (parámetros):}
\begin{itemize}
  \item PV: potencia nominal $\bar P^{\text{pv}}$, eficiencia/curva de disponibilidad.
  \item BESS: capacidad $E_{\max}$, límites $\bar P^{\text{ch}},\bar P^{\text{dis}}$, eficiencias $\eta_{\text{ch}},\eta_{\text{dis}}$, SoC mínimo/máximo.
  \item Red: precio $p_t$ (compra/venta), límite de intercambio $\bar P^{\text{grid}}$.
\end{itemize}

\end{frame}

\begin{frame}{Ejemplo de sistema 1: Microred con PV + batería + compra/venta a red}

\textbf{Estado (vector mínimo en $t$):}
\[
S_t=\big(\text{SoC}_t,\ \tilde G_t,\ \tilde L_t,\ p_t\big),
\]
donde $\tilde G_t$ (recurso solar) y $\tilde L_t$ (demanda) son entradas exógenas estocásticas (p.ej.\ Markov o AR).

\textbf{Eventos/actualización (modelo dinámico en pasos $\Delta t$):}
\[
\text{SoC}_{t+1}=\text{SoC}_t+\eta_{\text{ch}}P^{\text{ch}}_t\Delta t-\frac{1}{\eta_{\text{dis}}}P^{\text{dis}}_t\Delta t.
\]

\textbf{Decisiones y métricas:}
\[
a_t=(P^{\text{ch}}_t,P^{\text{dis}}_t,P^{\text{grid}}_t),\qquad
J=\mathbb E\Big[\sum_t p_t\,P^{\text{grid}}_t\Delta t+\lambda\,\text{ENS}_t\Big].
\]
\end{frame}

\begin{frame}{Ejemplo de sistema 2: Centro de distribución (fulfillment) con picking y packing}
\textbf{Propósito del estudio:} dimensionar recursos y política de asignación para reducir tiempo de ciclo y atrasos.

\medskip
\textbf{Entidades:}
\begin{itemize}
  \item \textbf{Órdenes} (jobs), \textbf{operarios} (pickers/packers), \textbf{estaciones} (picking/packing), \textbf{racks} y \textbf{transportador} (si aplica).
\end{itemize}

\textbf{Atributos:}
\begin{itemize}
  \item Orden $i$: tamaño $n_i$, prioridad $\pi_i$, fecha promesa $d_i$, zona(s) de picking.
  \item Operario $k$: velocidad $v_k$, habilidad, ubicación $loc_k(t)$.
  \item Estación $m$: capacidad (1 servidor), tasa de servicio dependiente del trabajo $T_m(\cdot)$.
\end{itemize}

\end{frame}


\begin{frame}{Ejemplo de sistema 2: Centro de distribución (fulfillment) con picking y packing}

\textbf{Estado (DES: simulación de eventos discretos):}
\[
S(t)=\Big(Q^{\text{pick}}(t),\ Q^{\text{pack}}(t),\ \{loc_k(t)\}_k,\ \{b_m(t)\}_m\Big),
\]
donde $Q^{\text{pick}},Q^{\text{pack}}$ son colas y $b_m(t)\in\{0,1\}$ indica si la estación $m$ está ocupada.

\textbf{Eventos (saltos en el estado):}
\begin{itemize}
  \item Llegada de orden, inicio/fin de picking, traslado entre zonas, inicio/fin de packing, despacho.
\end{itemize}

\textbf{Reglas de decisión (políticas):}
\[
a(t)=\text{asignación}( \text{orden}\rightarrow \text{operario/estación})\ \ \text{(p.ej.\ FIFO, prioridad, batching)}.
\]

\textbf{Métricas de desempeño:}
\[
\text{CycleTime}_i = C_i - A_i,\quad \text{Tardiness}_i=(C_i-d_i)_+,\quad
\mathbb E[\text{WIP}(t)],\ \mathbb P(\text{atraso}>0).
\]
\end{frame}



%================================================
% SLIDE 6: DETERMINÍSTICO VS ESTOCÁSTICO
%================================================
\begin{frame}{Sistemas: Determinísticos vs.\ Estocásticos}
\begin{columns}
\column{0.5\textwidth}
\textbf{Determinístico}
\begin{itemize}
  \item Sin aleatoriedad (o aleatoriedad despreciada).
  \item Misma entrada $\Rightarrow$ misma salida.
  \item Modelo: $Y = f(x)$.
  \item Útil si la variabilidad es irrelevante para la decisión.
\end{itemize}

\column{0.5\textwidth}
\textbf{Estocástico}
\begin{itemize}
  \item Al menos una variable aleatoria relevante.
  \item Salidas son variables aleatorias.
  \item Modelo: $Y = f(x,\xi)$, con $\xi$ aleatoria.
  \item Requiere inferencia estadística (promedios, intervalos, riesgo).
\end{itemize}
\end{columns}
\end{frame}

%================================================
% SLIDE 7: FORMALISMO MÍNIMO (PROBABILIDAD)
%================================================
\begin{frame}{Formalismo mínimo: ¿qué significa “aleatorio”?}
Trabajamos sobre un espacio de probabilidad $(\Omega,\mathcal F,\mathbb P)$.

\begin{itemize}
  \item Una \textbf{entrada aleatoria} es una función $\xi:\Omega\to\mathbb R^d$.
  \item Un \textbf{resultado de interés} es $Y = g(x, \xi)$ (p.ej.\ costo, utilidad, tiempo).
  \item En simulación, buscamos típicamente \textbf{esperanzas}:
  \[
  \mathbb E[Y] = \int_{\Omega} g(\xi(\omega))\, d\mathbb P(\omega).
  \]
\end{itemize}

\begin{block}{Idea central}
La simulación reemplaza integrales/esperanzas difíciles por promedios muestrales.
\end{block}
\end{frame}


%================================================
% SLIDE NUEVO A: EJEMPLO 1 (MUY BÁSICO) - MONEDA
%================================================
\begin{frame}{Ejemplo 1 (básico): lanzar una moneda}
\textbf{Experimento aleatorio:} lanzar una moneda una vez.

\medskip
\textbf{Espacio de resultados:}
\[
\Omega=\{C,S\}\quad \text{(Cara, Sello)}.
\]
\textbf{Evento:} $\{C\}\in\mathcal F$ (``sale cara''). \qquad
\textbf{Probabilidad:} $\mathbb P(C)=p,\ \mathbb P(S)=1-p$.

\medskip
\textbf{Entrada aleatoria (variable aleatoria):}
\[
\xi:\Omega\to\mathbb R,\qquad \xi(C)=1,\ \xi(S)=0.
\]
Interpretación: $\xi$ codifica el resultado del experimento en un número.

\medskip
\textbf{Resultado de interés:} $Y=g(\xi)$.
Por ejemplo, ``ganancia'':
\[
g(x)=2x-1\quad \Rightarrow\quad
Y=
\begin{cases}
+1,& \text{si sale cara}\\
-1,& \text{si sale sello}.
\end{cases}
\]

\end{frame}
\begin{frame}{Ejemplo 1 (básico): lanzar una moneda}

\medskip
\textbf{Esperanza (promedio teórico):}
\[
\mathbb E[Y]=\sum_{\omega\in\Omega} g(\xi(\omega))\,\mathbb P(\omega)
=(+1)\,p+(-1)(1-p)=2p-1.
\]

\begin{block}{Pregunta}
Si la expresión del valor esperado es muy díficil de cálcular que hacemos?.
\end{block}
\end{frame}

%================================================
% SLIDE NUEVO B: EJEMPLO 2 (BÁSICO) - DEMANDA Y COSTO
%================================================
\begin{frame}{Ejemplo 2:  demanda aleatoria y costo de inventario}
\textbf{Situación:} mañana la demanda $D$ puede ser baja o alta.

\medskip
\textbf{Espacio de resultados:}
\[
\Omega=\{\text{baja},\ \text{alta}\}.
\]
\textbf{Probabilidades:} $\mathbb P(\text{alta})=0.3$, $\mathbb P(\text{baja})=0.7$.

\medskip
\textbf{Entrada aleatoria:}
\[
\xi:\Omega\to\mathbb R,\qquad \xi(\text{baja})=80,\ \xi(\text{alta})=120.
\]
Aquí $\xi$ es literalmente la demanda aleatoria $D$.

\medskip
\textbf{Decisión fija:} ordenar $Q=100$ unidades.

\medskip
\textbf{Resultado de interés (costo por faltantes):}
\[
Y=g(\xi)=c_s\,( \xi - Q )_+,
\qquad (x)_+=\max\{x,0\},\qquad c_s=10.
\]
Interpretación: si la demanda supera $Q$, pagas $10$ por cada unidad que faltó; si no, el costo es 0.
\end{frame}

\begin{frame}{Ejemplo 2:  demanda aleatoria y costo de inventario}
\textbf{Situación:} mañana la demanda $D$ puede ser baja o alta.

\medskip
\textbf{Esperanza (promedio teórico):}
\[
\mathbb E[Y]=\sum_{\omega\in\Omega} g(\xi(\omega))\,\mathbb P(\omega)
=0\cdot 0.7 + \big(10(120-100)\big)\cdot 0.3 = 60.
\]

\begin{block}{Pregunta}
De que otra manera abordariamos esta situación?.
\end{block}
\end{frame}



%================================================
% SLIDE 8: JENSEN Y “FALACIA DEL PROMEDIO” (corregida)
%================================================
\begin{frame}{¿Por qué fallan los modelos “promedio”? (Jensen)}
\begin{alertblock}{Punto matemático}
Para $X$ aleatoria y $f$ no lineal, en general:
\[
\mathbb E[f(X)] \neq f(\mathbb E[X]).
\]
\end{alertblock}

\textbf{Jensen:}
\begin{itemize}
  \item Si $f$ es \textbf{convexa}: $\ \mathbb E[f(X)] \ge f(\mathbb E[X])$.
    \newline
    \emph{Lectura típica: costos/penalizaciones convexas} $\Rightarrow$ usar promedios \textbf{subestima} costo/riesgo.
  \item Si $f$ es \textbf{cóncava}: $\ \mathbb E[f(X)] \le f(\mathbb E[X])$.
    \newline
    \emph{Lectura típica: utilidades con saturación/capacidad} $\Rightarrow$ usar promedios \textbf{sobreestima} ganancia.
\end{itemize}
\end{frame}

\begin{frame}{Ejemplo concreto de Jensen (convexa y cóncava, con números)}
Sea $X$ una variable aleatoria con dos posibles valores:
\[
X=
\begin{cases}
0 & \text{con prob. } 0.5,\\
2 & \text{con prob. } 0.5.
\end{cases}
\qquad \Rightarrow \qquad \mathbb E[X]=1.
\]

\begin{columns}
\column{0.48\textwidth}
\textbf{Caso 1: $f$ convexa (costo)}\\
Tome $f(x)=x^2$.

\[
\mathbb E[f(X)]=\mathbb E[X^2]
=0.5(0^2)+0.5(2^2)=2.
\]
\[
f(\mathbb E[X])=(\mathbb E[X])^2=1^2=1.
\]

\begin{alertblock}{Conclusión}
\[
\mathbb E[X^2]=2 \;>\; 1=(\mathbb E[X])^2.
\]
\small
Si el ``costo'' crece de forma convexa (p.ej.\ penalización por retraso), usar el promedio ($X=1$) \textbf{subestima} el costo real esperado.
\end{alertblock}

\column{0.48\textwidth}
\textbf{Caso 2: $f$ cóncava (ganancia)}\\
Tome $f(x)=\sqrt{x}$.

\[
\mathbb E[f(X)]=\mathbb E[\sqrt{X}]
=0.5\sqrt{0}+0.5\sqrt{2}=0.707.
\]
\[
f(\mathbb E[X])=\sqrt{\mathbb E[X]}=\sqrt{1}=1.
\]

\begin{alertblock}{Conclusión}
\[
\mathbb E[\sqrt{X}]=0.707 \;<\; 1=\sqrt{\mathbb E[X]}.
\]
\small
Si la ``ganancia'' se satura (rendimientos decrecientes), usar el promedio \textbf{sobreestima} la ganancia esperada.
\end{alertblock}
\end{columns}

\end{frame}



%================================================
% SLIDE 9: DEFINICIÓN DE SIMULACIÓN
%================================================
\begin{frame}{Definición de Simulación}
\begin{quote}
\small
``La simulación es la imitación de la operación de un proceso o sistema del mundo real a lo largo del tiempo.'' --- Banks et al.
\end{quote}

\vspace{0.3cm}
Implica:
\begin{itemize}
  \item generar una \textbf{historia artificial} del sistema (trayectorias, escenarios);
  \item observar esa historia para estimar métricas (promedios, cuantiles, probabilidades);
  \item entender que el resultado es \textbf{numérico} y \textbf{estadístico} (con error).
\end{itemize}
\end{frame}



% --- Diapositiva con link + QR (lo más robusto en PDF) ---
\begin{frame}{Video de apoyo}
\centering
\Large\textbf{Video (YouTube): Simulación y aleatoriedad}

\vspace{0.5cm}
\href{https://www.youtube.com/watch?v=FCoLgj2ioJY}{\beamergotobutton{Abrir el video en YouTube}}

\vspace{0.8cm}
\qrcode[height=3.2cm]{https://www.youtube.com/watch?v=FCoLgj2ioJY}

\vspace{0.4cm}
\footnotesize\texttt{https://www.youtube.com/watch?v=FCoLgj2ioJY}
\end{frame}

%================================================
% SLIDE 10: CUÁNDO SÍ
%================================================
\begin{frame}{Cuándo es apropiada la simulación}
\begin{enumerate}
  \item El sistema es complejo (interacciones, restricciones, no linealidad).
  \item Hay incertidumbre relevante (variabilidad afecta la decisión).
  \item No hay solución analítica simple o ésta es frágil a supuestos.
  \item Se desea experimentar con diseños/políticas sin costo real.
  \item Se requiere estimar riesgo: $\mathbb P(Y>\tau)$, $\mathrm{VaR}$, cuantiles, colas.
\end{enumerate}
\end{frame}

%================================================
% SLIDE 11: CUÁNDO NO
%================================================
\begin{frame}{Cuándo NO es apropiada la simulación}
\begin{enumerate}
  \item El problema se resuelve con sentido común (o regla simple robusta).
  \item Hay solución analítica directa y suficiente (p.ej.\ EOQ simple bajo supuestos fuertes).
  \item Experimentar en el sistema real es más fácil/barato y seguro.
  \item No hay datos mínimos ni forma razonable de elicitar supuestos.
  \item El sistema no está bien definido: objetivos ambiguos, fronteras y métricas no claras.
\end{enumerate}
\end{frame}

%================================================
% SLIDE 12-13: EJEMPLO ALTERNATIVO (mejor que colas)
%================================================
\begin{frame}{Ejemplo mental: capacidad vs.\ variabilidad}
\textbf{Problema:} planear capacidad $C$ para cubrir demanda aleatoria $D$ (energía, inventario, producción).

\begin{itemize}
  \item Enfoque ``promedio'': fijar $C=\mathbb E[D]$ y concluir ``faltantes esperados = 0''.
  \item Realidad: el faltante es
  \[
    L(D) = (D-C)_+ = \max\{D-C,0\},
  \]
  que es una función \textbf{convexa} de $D$.
\end{itemize}

\begin{block}{Conclusión inmediata}
Aunque $C=\mathbb E[D]$, se tiene $\mathbb E[(D-C)_+] > 0$ si $\mathrm{Var}(D)>0$.
\end{block}
\end{frame}

\begin{frame}{Cálculo explícito (Normal): faltante esperado positivo}
Suponga $D\sim \mathcal N(\mu,\sigma^2)$ y planee $C=\mu$.

\[
\mathbb E[(D-\mu)_+] = \sigma\,\varphi(0) = \frac{\sigma}{\sqrt{2\pi}},
\]
donde $\varphi$ es la pdf Normal estándar.

\medskip
Ejemplo numérico: $\mu=100$, $\sigma=30$:
\[
\mathbb E[(D-100)_+] \approx \frac{30}{\sqrt{2\pi}} \approx 11.97.
\]

\textbf{Lectura:} ``en promedio'' faltan \(\approx 12\) unidades por periodo, aun cuando la capacidad iguala la media.  
Esto es \textbf{variabilidad pura} + \textbf{convexidad} (Jensen).
\end{frame}

%================================================
% SLIDE 14-15: PROS/CONS
%================================================
\begin{frame}{Ventajas de la simulación}
\begin{itemize}
  \item \textbf{Control del tiempo:} compresión/expansión temporal.
  \item \textbf{Control experimental:} comparar políticas con mismas realizaciones aleatorias (CRN).
  \item \textbf{Seguridad:} experimentar sin intervenir el sistema real.
  \item \textbf{Riqueza de métricas:} no solo medias, también cuantiles, colas y probabilidades.
\end{itemize}
\end{frame}

\begin{frame}{Desventajas de la simulación}
\begin{itemize}
  \item Requiere habilidad en modelado, probabilidad y programación.
  \item Resultados son \textbf{estimaciones} (con error); interpretar requiere estadística.
  \item Construir un modelo \textbf{válido} toma tiempo (y datos).
  \item Riesgo de \emph{``Garbage In, Garbage Out''}: supuestos malos producen conclusiones falsas.
\end{itemize}
\end{frame}




%================================================
% SLIDE 16: TIPOS 
%================================================

%================================================
% 3 DIAPOSITIVAS: UNIVERSO DE LA SIMULACIÓN + DÓNDE ENTRA MONTE CARLO
%================================================

\begin{frame}{El universo de la simulación: ¿qué significa “tipo” de simulación?}
Un \textbf{tipo de simulación} nace de \textbf{cómo representamos la evolución del sistema}.

\vspace{0.3cm}
\begin{block}{Decisión 1: ¿hay evolución temporal?}
\begin{itemize}
  \item \textbf{Simulación estática (sin tiempo):} el sistema se estudia en un “instante/escenario”.
  \item \textbf{Simulación dinámica (con tiempo):} existe un estado $S(t)$ que cambia.
\end{itemize}
\end{block}

\begin{block}{Decisión 2 (si es dinámica): ¿cómo cambia el estado?}
\begin{itemize}
  \item \textbf{Eventos discretos (DES):} $S(t)$ cambia solo cuando ocurre un evento (saltos).
  \item \textbf{Por pasos de tiempo:} $S_{t+1}=F(S_t,\ldots)$ en periodos ($t=0,1,2,\dots$).
  \item \textbf{Continuo:} $S(t)$ cambia continuamente (ecuaciones diferenciales).
  \item \textbf{Basada en agentes / híbrida:} reglas locales (agentes) y/o mezcla de mecanismos.
\end{itemize}
\end{block}

\begin{alertblock}{Idea clave}
DES/continuo/pasos/agentes describen \textbf{la mecánica del modelo} (cómo evoluciona el sistema).
\end{alertblock}
\end{frame}


\begin{frame}{¿De dónde salen los tipos? (se crean al definir estado + tiempo)}
Para construir un simulador siempre se hacen estas elecciones:

\vspace{0.2cm}
\begin{enumerate}
  \item \textbf{Definir el estado} (variables mínimas para describir el sistema):
  \[
  S(t)=(x_1(t),\dots,x_n(t)).
  \]
  \item \textbf{Definir la regla de evolución (el “reloj” del modelo):}
  \begin{itemize}
    \item \textbf{DES:} el tiempo avanza al \emph{siguiente evento} (llegada, fin de servicio, falla).
    \item \textbf{Pasos:} el tiempo avanza en incrementos fijos $\Delta t$ (hora a hora, día a día).
    \item \textbf{Continuo:} el tiempo es continuo y $S(t)$ sigue una dinámica (ODE/SDE).
  \end{itemize}
  \item \textbf{Definir qué es aleatorio} (entradas inciertas):
  \[
  \xi \ \text{(demanda, tiempos, fallas, clima, comportamiento)}.
  \]
  \item \textbf{Definir qué se mide} (salida de interés):
  \[
  Y = g(\text{modelo},\xi)\quad \text{(costo, tiempo, utilidad, riesgo)}.
  \]
\end{enumerate}

\begin{block}{Lectura}
El \textbf{tipo} queda determinado por (estado + forma de avanzar el tiempo).  
La \textbf{incertidumbre} se incorpora en $\xi$.
\end{block}
\end{frame}


\begin{frame}{¿Dónde entra Monte Carlo ?}
Una vez definido el modelo (estático o dinámico), suele interesarnos una métrica \textbf{promedio o de riesgo}.

\vspace{0.2cm}
\begin{block}{Preguntas típicas}
\[
\mathbb E[Y]\quad (\text{costo/tiempo/utilidad promedio}),\qquad
\mathbb P(Y>\tau)\quad (\text{riesgo}),\qquad
q_\alpha(Y)\quad (\text{percentiles}).
\]
\end{block}

\begin{alertblock}{Monte Carlo}
\textbf{Monte Carlo} es el procedimiento de:
\begin{itemize}
  \item generar muchas realizaciones de la incertidumbre ($\xi^{(1)},\dots,\xi^{(N)}$),
  \item obtener la salida en cada una ($Y^{(i)}$),
  \item y promediar/contar para estimar lo que queremos.
\end{itemize}
\[
\mathbb E[Y] \approx \frac{1}{N}\sum_{i=1}^N Y^{(i)}.
\]
\end{alertblock}
\end{frame}


\begin{frame}{¿Que veremos en el curso respecto a la simulación?}
\begin{block}{Lo que haremos en el curso}
\begin{itemize}
  \item \textbf{Simulación Monte Carlo} (escenarios/réplicas para estimar promedios y riesgo).
  \item \textbf{Simulación de eventos discretos (DES)} (modelos dinámicos cuyo estado cambia en eventos).
\end{itemize}
\end{block}
\end{frame}




%================================================
% SLIDE 17-22: NEWSVENDOR (corregido y reforzado)
%================================================
\begin{frame}{Caso de estudio: Newsvendor (vendedor de periódicos)}
\textbf{Decisión:} escoger cantidad $Q$ de producto perecedero para vender mañana.

\begin{itemize}
  \item Costo por unidad: $c=20$.
  \item Precio: $p=50$.
  \item Salvamento: $s=0$.
\end{itemize}

\textbf{Demanda:} $D$ aleatoria. Suponga:
\[
D \sim \mathcal N(\mu=100,\sigma=30).
\]

\medskip
\textbf{Utilidad:}
\[
\Pi(D,Q) = p\min(D,Q) + s(Q-\min(D,Q)) - cQ
= (p-s)\min(D,Q) + (s-c)Q.
\]
\end{frame}

\begin{frame}{Enfoque ingenuo: fijar demanda a la media}
\textbf{Heurística ``promedio'':} asumir $D=\mu$ y pedir $Q=\mu=100$.

\[
\Pi(\mu,\mu) = (p-c)\mu = (50-20)\cdot 100 = 3000.
\]

\begin{alertblock}{Problema}
Ese número no es $\mathbb E[\Pi(D,Q)]$. Es la utilidad de un mundo sin incertidumbre (o con información perfecta).
\end{alertblock}
\end{frame}

\begin{frame}{Por qué la esperanza baja: concavidad en la demanda}
Para $Q$ fijo, $\Pi(D,Q)$ como función de $D$ es:
\[
\Pi(D,Q)=p\min(D,Q)-cQ,
\]
una función \textbf{cóncava} en $D$ (pendiente $p$ hasta $Q$ y luego pendiente $0$).

\begin{block}{Jensen aplicado}
Como $\Pi(\cdot,Q)$ es cóncava:
\[
\mathbb E[\Pi(D,Q)] \le \Pi(\mathbb E[D],Q).
\]
\end{block}

\textbf{Lectura:} el ``promedio'' tiende a \textbf{sobreestimar} ganancias cuando hay saturación/capacidad.
\end{frame}

\begin{frame}{Solución analítica (óptimo estocástico)}
En newsvendor clásico:
\[
F(Q^\star) = \frac{p-c}{p-s}
= \frac{50-20}{50-0}=0.6.
\]

Si $D\sim \mathcal N(\mu,\sigma^2)$,
\[
Q^\star = \mu + \sigma z_{0.6},
\quad z_{0.6}\approx 0.2533
\Rightarrow
Q^\star \approx 100 + 30(0.2533) \approx 107.6.
\]

\textbf{Interpretación:} se ordena por encima de la media porque el costo de escasez (margen perdido) supera el costo de exceso.
\end{frame}

\begin{frame}{Cálculo analítico rápido: $\mathbb E[\min(D,Q)]$ (Normal)}
Sea $D\sim \mathcal N(\mu,\sigma^2)$ y defina $z=\frac{Q-\mu}{\sigma}$.
Entonces:
\[
\mathbb E[\min(D,Q)]
= \mu \Phi(z) - \sigma \varphi(z) + Q\bigl(1-\Phi(z)\bigr),
\]
donde $\Phi$ y $\varphi$ son la cdf/pdf de la Normal estándar.

\medskip
Con $\mu=100$, $\sigma=30$, $Q=100$ ($z=0$):
\[
\mathbb E[\min(D,100)]
=100 - 30\varphi(0)
\approx 100 - 30(0.3989)\approx 88.03.
\]

Por tanto:
\[
\mathbb E[\Pi(D,100)] = p\,\mathbb E[\min(D,100)] - c\cdot 100
\approx 50(88.03) - 2000 \approx 2401.6.
\]
\end{frame}

\begin{frame}{Métricas estándar: EEV, RP, VSS, EVPI}
Definiciones (clásicas en programación estocástica):

\begin{itemize}
  \item \textbf{EEV:} evaluar la solución del ``valor esperado'' (aquí $Q=\mu=100$):
  \[
  \text{EEV}=\mathbb E[\Pi(D,100)]\approx 2401.6.
  \]
  \item \textbf{RP:} valor óptimo estocástico (aquí $Q^\star\approx107.6$):
  \[
  \text{RP}=\mathbb E[\Pi(D,Q^\star)]\approx 2420.5.
  \]
  \item \textbf{VSS:} valor de resolver el problema estocástico:
  \[
  \text{VSS}=\text{RP}-\text{EEV}\approx 18.9.
  \]
  \item \textbf{EVPI:} valor de información perfecta:
  \[
  \text{EVPI}=\mathbb E\big[\max_Q \Pi(D,Q)\big]-\text{RP}.
  \]
  Con información perfecta y $s=0$, $\max_Q \Pi(D,Q)=(p-c)D$, así
  $\mathbb E[(p-c)D]=3000$ y $\text{EVPI}\approx 3000-2420.5\approx 579.5$.
\end{itemize}
\end{frame}

%================================================
% SLIDE 23: DIAGRAMA MC
%================================================
\begin{frame}{Diagrama de flujo (Monte Carlo)}
Para $i=1,\dots,N$:
\begin{enumerate}
  \item Generar $d_i \sim \mathcal N(100,30)$ (o variante factible no-negativa).
  \item Calcular ventas $v_i=\min(d_i,Q)$.
  \item Calcular utilidad $\pi_i = p v_i - cQ$.
\end{enumerate}

Estimador:
\[
\widehat{\mathbb E[\Pi]}=\bar \pi_N=\frac{1}{N}\sum_{i=1}^N \pi_i.
\]

\begin{block}{Error estadístico (CLT)}
\[
\bar \pi_N \pm z_{0.975}\frac{s_\pi}{\sqrt{N}}
\]
da un intervalo de confianza aproximado al 95\%.
\end{block}
\end{frame}

%================================================
% SLIDE 24-25:
%================================================
\begin{frame}[fragile]{Implementación en Python: simulación + chequeo analítico}
\end{frame}

\begin{frame}[fragile]{Implementación en Python: simulación + chequeo analítico}
\end{frame}

\begin{frame}[fragile]{Implementación en Python: estimación + intervalo de confianza}
\end{frame}

%================================================
% SLIDE 26: RESULTADOS (corregidos)
%================================================
\begin{frame}{Resultados}
Con los parámetros $\mu=100,\sigma=30,Q=100,p=50,c=20$:

\begin{itemize}
  \item \textbf{Benchmark sin incertidumbre:} $\Pi(\mu,\mu)=(p-c)\mu=3000$.
  \item \textbf{Esperanza analítica:} $\mathbb E[\Pi(D,100)]\approx 2401.6$.
  \item \textbf{Monte Carlo típico (semilla 42, $N=10^4$):}
  \[
  \widehat{\mathbb E[\Pi]}\approx 2399.0
  \quad \text{con IC95\% } \pm 17.2.
  \]
  \item \textbf{Pérdida:} $\mathbb P(\Pi<0)=\mathbb P(D<40)\approx \Phi(-2)\approx 0.0228$ (coincide con MC $\approx 0.023$).
\end{itemize}

\begin{alertblock}{Mensaje}
El promedio ``optimista'' (3000) ignora saturación y riesgo a la baja; la ganancia esperada real cae cerca de \(\sim 2400\) (brecha $\sim 600$, \(\sim 20\%\)).
\end{alertblock}
\end{frame}

%================================================
% SLIDE 27: VISUALIZACIÓN (igual, pero texto ajustado)
%================================================
\begin{frame}[fragile]{Visualización de la distribución de utilidades}
\end{frame}

%================================================
% SLIDE 28: INTERPRETACIÓN
%================================================
\begin{frame}{Interpretación (qué mirar en el histograma)}
\begin{itemize}
  \item \textbf{Truncamiento a la derecha:} para $D>Q$ las ventas se saturan en $Q$ y la utilidad se ``aplana''.
  \item \textbf{Cola izquierda:} eventos de demanda baja generan utilidades mucho menores; esa asimetría reduce la media.
  \item \textbf{Error Monte Carlo:} el ancho del IC decrece como $O(1/\sqrt{N})$; pedir más precisión requiere muchas más corridas.
\end{itemize}

\textbf{Preguntas guía:}
\begin{itemize}
  \item ¿Qué pasa con la media si $\sigma$ sube a 50?
  \item ¿Cómo cambia $Q^\star$ cuando cambia $(p-c)$ o $(c-s)$?
\end{itemize}
\end{frame}

%================================================
% SLIDE 29-31: EJERCICIOS
%================================================
%================================================
% EJERCICIOS ESCALADOS (reemplaza tus slides 26-28)
%================================================

\begin{frame}{Ejercicio: ¿Dónde está la incertidumbre? (identificar variables aleatorias)}
En cada situación, responda \textbf{tres cosas}:
\[
\text{(i) ¿Qué es aleatorio?}\quad
\text{(ii) ¿Qué se observa?}\quad
\text{(iii) ¿Qué métrica importa?}
\]

\vspace{0.2cm}
\begin{enumerate}
  \item \textbf{Cafetería:} llegan estudiantes y hay 1 caja. A veces se forma fila.
  \item \textbf{Inventario:} una tienda decide $Q$ unidades para mañana; no sabe cuántas venderá.
  \item \textbf{Transporte:} un mensajero tarda distinto según tráfico.
\end{enumerate}

\end{frame}




\begin{frame}{Ejercicio: Modelar la incertidumbre (Casos A--D)}
\small
\begin{columns}[T]
\column{0.45\textwidth}

\begin{block}{Caso A}
En una ventanilla, atender a un cliente siempre requiere completar \textbf{tres pasos consecutivos} (verificar datos, registrar, confirmar pago).  
Cada paso toma tiempo positivo y el tiempo total observado suele concentrarse alrededor de un valor típico, con variabilidad moderada.
\end{block}

\begin{block}{Caso B}
En un punto de información, los visitantes llegan sin coordinación.  
A veces llegan varios casi seguidos; otras veces hay pausas. El tiempo entre llegadas es siempre positivo y no se observa un “máximo” claro.
\end{block}

\column{0.45\textwidth}

\begin{block}{Caso C}
Durante cada hora, un call center registra un \textbf{número de llamadas entrantes}.  
Se trata de un conteo (0,1,2,...) y el equipo quiere modelar cuántas llamadas llegan por hora para dimensionar turnos.
\end{block}

\begin{block}{Caso D}
Una empacadora llena bolsas con una máquina calibrada a 500 g.  
El peso final varía ligeramente por vibraciones y tolerancias: algunas bolsas quedan por encima y otras por debajo; la mayoría queda cerca del objetivo y los extremos son poco frecuentes.
\end{block}

\end{columns}
\end{frame}


\begin{frame}{Ejercicio: Modelar la incertidumbre (Casos E--H)}
\small
\begin{columns}[T]
\column{0.45\textwidth}

\begin{block}{Caso E}
El tiempo de reparación de un equipo suele ser corto cuando el diagnóstico es directo,  
pero ocasionalmente se alarga mucho por espera de repuesto, recalibración o retrabajo.  
La variable es positiva y se observan pocos valores muy grandes.
\end{block}

\begin{block}{Caso F}
Se estudia la vida útil de un componente que se degrada con el uso.  
Al inicio casi no falla, pero a medida que envejece la probabilidad de falla aumenta.  
Se registra el tiempo hasta falla (positivo) para planear mantenimiento.
\end{block}

\column{0.45\textwidth}

\begin{block}{Caso G}
Un bus interno recorre siempre la misma ruta dentro del campus.  
Por restricciones de distancia, velocidad y semáforos, el tiempo de viaje está siempre entre 12 y 18 minutos, y no hay evidencia clara de que algún valor dentro del rango sea “especialmente más probable”.
\end{block}

\begin{block}{Caso H}
Para una actividad de proyecto (instalar un módulo), un experto da tres estimaciones:  
\textbf{optimista} (si todo sale perfecto), \textbf{más probable} (lo usual) y \textbf{pesimista} (si aparece un problema común).  
Se quiere modelar la duración con esa información.
\end{block}

\end{columns}
\end{frame}



\begin{frame}{Ejercicio: distinguir \textit{variables} vs \textit{parámetros} (y qué es aleatorio)}
\textbf{Objetivo:} en cada situación, clasifique cada símbolo en una de estas categorías:
\[
\textbf{(V)}\ \text{variable de decisión/estado} \qquad
\textbf{(PA)}\ \text{parámetro aleatorio (variable aleatoria)} \qquad
\textbf{(PD)}\ \text{parámetro determinístico (constante)}
\]

\begin{block}{Situación 1: Inventario (un solo día)}
Una tienda decide $Q$ unidades para mañana. La demanda del día es $D$.  
Precio $p=50$, costo $c=20$, salvamento $s=0$.  
Utilidad: $\Pi(D,Q)=p\min(D,Q)+s(Q-\min(D,Q)) - cQ$.

\textbf{Clasifique:} $Q,\ D,\ p,\ c,\ s,\ \Pi(D,Q)$.
\end{block}


\begin{block}{Situación 2: Caja de supermercado (servicio)}
Llegan clientes. El tiempo entre llegadas es $A$ (min) y el tiempo de servicio es $S$ (min).  
El sistema tiene una sola caja.

\textbf{Clasifique:} $A,\ S,\ \lambda,\ \mu,\ \rho,\ N(t)$, donde:
\begin{itemize}
  \item $\lambda$ = tasa promedio de llegadas (clientes/min),
  \item $\mu$ = tasa promedio de servicio (clientes/min),
  \item $\rho=\lambda/\mu$ = utilización,
  \item $N(t)$ = número de clientes en el sistema en el tiempo $t$.
\end{itemize}
\end{block}

\end{frame}


%================================================
% SLIDE 33: REFERENCIAS
%================================================
\begin{frame}{Referencias bibliográficas}
\begin{itemize}
  \item Banks, J., Carson, J. S., Nelson, B. L., \& Nicol, D. M. (2010). \textit{Discrete-Event System Simulation}. Prentice Hall.
  \item Law, A. M. (2015). \textit{Simulation Modeling and Analysis}. McGraw-Hill.
  \item Kulkarni, V. G. (2011). \textit{Introduction to Modeling and Analysis of Stochastic Systems}. Springer.
  \item Ross, S. M. (2013). \textit{Simulation}. Academic Press. (Monte Carlo, CLT, intervalos)
\end{itemize}
\end{frame}

\end{document}


\begin{frame}{Ejercicio 2: Modelar la incertidumbre (sin “adivinar” distribuciones)}
Para cada caso elija una forma simple de modelar la incertidumbre y \textbf{justifique con 2 razones}.

\vspace{0.2cm}
\begin{columns}[T]
\column{0.5\textwidth}
\textbf{Caso A: tiempo de servicio en caja} (minutos)\\
Opciones:
\begin{itemize}
  \item \textbf{Exponencial} (si “no hay memoria”, alta variabilidad).
  \item \textbf{Gamma/Erlang} (si son varias etapas).
  \item \textbf{Lognormal} (si es positivo y muy asimétrico).
  \item \textbf{Empírico} (si hay datos: muestrear de la lista).
\end{itemize}

\column{0.5\textwidth}
\textbf{Caso B: demanda diaria} (unidades)\\
Opciones:
\begin{itemize}
  \item \textbf{Normal truncada} (si es alrededor de una media, pero no negativa).
  \item \textbf{Poisson} (si son conteos de “eventos”).
  \item \textbf{Empírico / bootstrap} (si hay histórico).
\end{itemize}
\end{columns}

\vspace{0.2cm}
\begin{alertblock}{Regla de oro}
Antes de elegir distribución, revise: \textbf{soporte} (¿puede ser negativo?), \textbf{simetría}, \textbf{colas}, y si el fenómeno es \textbf{discreto} o \textbf{contínuo}.
\end{alertblock}
\end{frame}



\begin{frame}[fragile]{Ejercicio 3: Primer Monte Carlo en Python (llenar huecos)}
Objetivo: estimar $\mathbb E[\text{costo}]$ cuando la demanda es incierta.

\vspace{0.2cm}
\textbf{Modelo:} Usted ordena $Q=100$. Si la demanda $D$ supera $Q$, paga un costo de faltante $c_s=10$ por unidad faltante:
\[
Y = c_s (D-Q)_+.
\]


\begin{block}{Preguntas}
(i) ¿Qué cambia si $N$ pasa de 5{,}000 a 50{,}000? \quad
(ii) ¿Qué cambia si sube $\sigma$?
\end{block}
\end{frame}


\begin{frame}{Ejercicio 4: Idea de DES sin pseudocódigo (simulación manual por eventos)}
\textbf{Sistema:} una caja, una fila. Llegan clientes y la caja atiende uno a la vez.

\vspace{0.2cm}
\begin{block}{Estado mínimo}
\[
S(t)=\big(\text{\# en fila},\ \text{caja libre/ocupada},\ \text{tiempo de fin de servicio (si está ocupada)}\big).
\]
\end{block}

\textbf{Eventos posibles:}
\begin{itemize}
  \item \textbf{Llegada:} llega un cliente.
  \item \textbf{Fin de servicio:} termina el cliente actual.
\end{itemize}

\vspace{0.2cm}
\textbf{Actividad (en grupos):} complete a mano los \textbf{primeros 6 eventos} en una tabla.
\begin{center}
\small
\begin{tabular}{@{}cccccc@{}}
\toprule
Evento \# & Tiempo & Tipo & \# en fila & Caja (L/O) & Próx.\ fin servicio \\
\midrule
1 & 0.0 & Llegada &  &  &  \\
2 & 1.2 & Llegada &  &  &  \\
3 & 2.0 & Fin serv. &  &  &  \\
4 & 2.4 & Llegada &  &  &  \\
5 & 3.1 & Llegada &  &  &  \\
6 & 3.7 & Fin serv. &  &  &  \\
\bottomrule
\end{tabular}
\end{center}

\begin{alertblock}{Regla (sin código)}
En \textbf{llegada}: si la caja está libre $\Rightarrow$ inicia servicio; si no $\Rightarrow$ aumenta fila.\\
En \textbf{fin de servicio}: si hay fila $\Rightarrow$ pasa el siguiente; si no $\Rightarrow$ la caja queda libre.
\end{alertblock}
\end{frame}

\end{document}
